\section{Disclosed Limitation: TLS Certificate Validation}
\label{sec:tls}

\subsection{The Limitation}

\bisect{} uses HTTPS for all Census Bureau downloads.
TLS certificate validation is performed against the system trust store
(the operating system's certificate authority database).
\bisect{} does \emph{not} implement certificate pinning---it does not
hard-code the specific TLS certificate or public key of the Census Bureau's
server.

This creates a potential attack surface for man-in-the-middle (MITM)
interception:
\begin{enumerate}
\item An adversary who can add a certificate authority to the system trust
  store (requires administrator access to the redistricting machine) could
  intercept TLS connections to \texttt{www2.census.gov} and serve
  manipulated data.
\item The manipulated data would be downloaded with a valid TLS certificate
  (from the adversary's CA) and would not trigger a TLS validation error.
\item The \bisect{} SHA-256 hash would record the hash of the \emph{manipulated}
  file, not the original.
  This hash would not match the Census Bureau's published checksum,
  but the adversary controls the comparison too (if they also control
  the manifest).
\end{enumerate}

This attack chain requires the adversary to simultaneously:
(a) compromise the system trust store;
(b) intercept network traffic to the Census Bureau;
(c) serve manipulated data that is plausible enough to pass visual inspection;
(d) modify the comparison checksum source.

\subsection{Risk Assessment}

This attack requires infrastructure-level compromise of the redistricting
authority's computing environment---comparable to compromising the system
running the redistricting algorithm entirely.
An adversary with this level of access could also directly modify the
algorithm binary, the parameter files, or the output plan, bypassing the
data layer entirely.

In the context of the \dia{}'s deployment model (redistricting run by
a certified neutral authority with audited computing infrastructure),
this attack is not plausible without a broader security failure.

\subsection{Recommended Mitigation}

For litigation-grade deployments, we recommend supplementary out-of-band
hash verification:
\begin{enumerate}
\item Download Census P.L.\ 94-171 files via two independent network paths
  (e.g., direct download and Census FTP mirror).
\item Compare SHA-256 hashes between the two copies.
\item Additionally compare against hashes published by independent
  redistricting data services (e.g., IPUMS NHGIS, Harvard Election Data Archive).
\end{enumerate}
This multi-source verification eliminates the TLS limitation without
requiring certificate pinning.

Future versions of \bisect{} may implement certificate pinning using the
Census Bureau's published TLSA DNS records (RFC 6698), which would
eliminate the attack surface entirely.
This is tracked as a known limitation in the \bisect{} security documentation.
